\subsection*{O228\quad 正弦相位光栅的Bessel衍射级、能量守恒与零级抑制}
单色相干平面光正入射无损薄相位光栅，周期$d$，复透过率$t(x)=\exp[i m\cos(2\pi x/d)]$，相位调制深度$m\ge0$，忽略传播中的偏振和材料吸收。
(1) 用Jacobi–Anger展开推导离散衍射级复振幅。
(2) 给出各级次相对能量及传播角条件。
(3) 在弱调制极限推导零级与$\pm1$级强度分配并检验能量守恒。
(4) 求零级被完全抑制的第一个调制深度并讨论更高级次能量。
\par\textbf{逐问解析}
(1) 周期相位因子的傅里叶展开$\exp(im\cos\theta)=\sum_{n=-\infty}^\infty i^nJ_n(m)e^{in\theta}$。因此透射波可分解为横向波数$k_{x,n}=2\pi n/d$的平面波，每级复幅$a_n=i^nJ_n(m)$；该结果来自振幅的相干傅里叶分解，而不是对入射强度直接调制。
(2) 若入射、出射介质折射率均为1，则级次$n$满足$\sin\theta_n=n\lambda/d$，仅$|n\lambda/d|\le1$的级次为远场传播模；在所有衍射级均可传播、并适当使用功率归一化的理想近轴模型中，级次功率分数为$P_n/P_{\rm in}=J_n(m)^2$，且$\sum_{n=-\infty}^\infty J_n(m)^2=1$。若部分级次倏逝，不能把所有$J_n^2$无修正地解释为远场输送功率。
(3) 当$m\ll1$，$J_0(m)=1-m^2/4+O(m^4)$、$J_1(m)=m/2+O(m^3)$，故零级强度$J_0^2=1-m^2/2+O(m^4)$，一级各自$J_{\pm1}^2=m^2/4+O(m^4)$；三项合计$1+O(m^4)$。在二阶弱调制下，零级损失的能量恰转入对称的一级衍射，不凭空产生能量。
(4) 要消除零级，须$J_0(m)=0$，其首个正根$m_1\simeq2.404826$。此时大约$J_{\pm1}(m_1)^2\simeq0.2695$分别进入$\pm1$级，两者总和仅约0.539，余下能量进入$\pm2,\pm3,\ldots$高阶；因此零级完全抑制不意味着一级级次总效率为100%。有限空间口径和光栅色散还会改变实际角谱。
\par\medskip
\subsection*{O229\quad 单总线微环共振器的临界耦合与吸收线宽}
单总线侧耦合光学微环采用时间耦合模理论，腔本征角频率$\omega_0$，内损耗振幅衰减率$\gamma_i\ge0$、总线外耦合振幅衰减率$\gamma_e>0$，输入输出波振幅满足功率归一化$|s|^2=P$。
(1) 建立共振腔振幅方程及输入—输出关系。
(2) 求通带幅透射系数与功率透射谱。
(3) 推导临界耦合条件与总功率去向。
(4) 求谱线全宽、加载品质因数，并讨论无内损极限。
\par\textbf{逐问解析}
(1) 取场的时间因子$e^{-i\omega t}$，定义$\dot a=(-i\omega_0-\gamma_i-\gamma_e)a+\sqrt{2\gamma_e}s_{\rm in}$与$s_{\rm out}=s_{\rm in}-\sqrt{2\gamma_e}a$，使$|a|^2$为腔内储能。自由无驱动时$d|a|^2/dt=-2(\gamma_i+\gamma_e)|a|^2$，所以这里$\gamma$表示振幅衰减率，而非能量衰减率。
(2) 令$\Delta=\omega-\omega_0$、$a(t)=a_\omega e^{-i\omega t}$，则$a_\omega=\sqrt{2\gamma_e}s_{\rm in}/(\gamma_i+\gamma_e-i\Delta)$。由输入—输出式有$t(\Delta)=s_{\rm out}/s_{\rm in}=[\gamma_i-\gamma_e-i\Delta]/[\gamma_i+\gamma_e-i\Delta]$，故$T(\Delta)=|t|^2=[(\gamma_i-\gamma_e)^2+\Delta^2]/[(\gamma_i+\gamma_e)^2+\Delta^2]$。
(3) 共振时若$\gamma_i=\gamma_e$，则$t(0)=0$，总线直通端完全熄灭，这就是临界耦合。由$1-T=4\gamma_i\gamma_e/[(\gamma_i+\gamma_e)^2+\Delta^2]$可得共振被吸收的功率分数，临界耦合时恰为1；若$\gamma_i=0$则不可能有真正能量吸收，即使共振会产生强相位变化。
(4) 透射谱的缺口$1-T$是宽度参数$\Gamma=\gamma_i+\gamma_e$的Lorentz函数，在$|\Delta|=\Gamma$时降为峰值的一半，所以全宽$\Delta\omega_{\rm FWHM}=2(\gamma_i+\gamma_e)$。加载品质因数$Q_L=\omega_0/[2(\gamma_i+\gamma_e)]$；无内损$\gamma_i\to0$时功率透射$T\to1$对所有频率成立，但相位仍在共振附近急变，不能把全通率误判为没有腔内储能。
\par\medskip
